%% Bilingual Homework Template starter -- homework.cls by Yuzheng Cai,
%% itself built on Josh Davis's well-known latex-homework-template
%% (credited in the class and in references.bib).
%%
%% homework.cls IS BUNDLED here (not left to the system) because a
%% DIFFERENT, unrelated CTAN class -- Jinwen Xu's 2023-2025
%% homework.cls -- already ships in every modern TeX distribution
%% under the SAME filename (the same collision hit math-homework
%% elsewhere in this batch). A bare \documentclass{homework} on a
%% stock TeX Live install silently resolves to that unrelated class
%% instead -- confirmed by probing this pipeline's own pinned image.
%%
%% Source: https://github.com/YZ-Cai/LatexHomeworkTemplate
%% Licence: MIT, (c) 2022 Yuzheng Cai.
%%
%% engine.txt: xelatex (required -- the class \usepackage{ctex}
%% unconditionally, for xeCJK, regardless of the E/C header-language
%% toggle below).
%%
%% This class's chrome language (cover page, running header/footer --
%% "English"/"Chinese") is chosen by editing homework.cls itself
%% (line 5 vs 6), not by a \documentclass option. Left at the class's
%% own default (English chrome) here: this class is explicitly
%% documented as "English template, but ALSO SUPPORTS CHINESE IN THE
%% DOCUMENT BODY" -- i.e. genuinely bilingual within one document,
%% which this starter demonstrates directly (English problem
%% statements, Chinese worked solutions) rather than switching all
%% chrome to Chinese.
\documentclass{homework}

\newcommand{\hmwkTitle}{Homework \#4}
\newcommand{\hmwkDueDate}{January 8, 2027}
\newcommand{\hmwkClass}{Natural Language Processing}
\newcommand{\hmwkAuthorName}{Ada Lovelace}

\begin{document}
\maketitle
\if\hmwkCoverPage 1
    \pagebreak
\fi

\begin{homeworkProblem}
Explain why byte-pair encoding (BPE) tokenization can represent
out-of-vocabulary words that a fixed word-level vocabulary cannot.

\solution
BPE 从字符级别开始，通过统计训练语料中相邻符号对的出现频率，反复合并最高频的符号对，
逐步构建一个包含子词单元的词表。由于任意单词都可以被拆分为若干个已知的子词片段
（在最坏情况下可以退化到单个字符），因此即使某个完整单词从未在训练语料中出现过，
BPE 分词器仍然可以用已知的子词单元将其表示出来，而不会产生“未登录词”（OOV）问题。
\end{homeworkProblem}

\begin{homeworkProblem}
Given the bigram counts $C(\text{the}, \text{cat}) = 12$ and
$C(\text{the}) = 400$, compute the maximum-likelihood estimate of
$P(\text{cat} \mid \text{the})$, and explain one drawback of using raw
MLE counts for language modeling.

\solution
The maximum-likelihood bigram probability is
\begin{equation}
    P(\text{cat} \mid \text{the}) = \frac{C(\text{the}, \text{cat})}{C(\text{the})}
    = \frac{12}{400} = 0.03.
\end{equation}
一个主要缺陷是：任何在训练语料中从未共现过的二元组都会被赋予概率 0，
这会导致模型在遇到测试集中真实存在但训练时未见过的词对时，将整句话的概率计算为 0
（数据稀疏问题）。常见的缓解方法包括加一平滑（Laplace smoothing）和 Kneser-Ney 平滑等。
\end{homeworkProblem}

\begin{homeworkProblem}
Try citation. For example, thanks Josh Davis for his original work on
the latex-homework-template project \cite{LatexHomeworkTemplate}.

\solution
Citations use plain \verb|bibtex| via \verb|\bibliographystyle{abbrv}|
and \verb|\bibliography{references.bib}| at the end of the document.
\end{homeworkProblem}

\bibliographystyle{abbrv}
\bibliography{references.bib}

\end{document}
